\section{Ортогональные системы функций}\label{sec-fourier-orthogonal}

Здесь рассматриваются вещественные функции на отрезке \([a,b]\), \(a<b\). Интегралы понимаются по Риману; далее также допускаются абсолютно сходящиеся несобственные интегралы с конечным числом особых точек.

\begin{definition}[Скалярное произведение и норма]\label{def-function-inner-product}

Для \(f,g\in C([a,b])\) положим
\begin{equation}
(f,g)=\int_a^b f(x)g(x)\,dx,
\qquad
\|f\|=\sqrt{(f,f)}=\left(\int_a^b f^2(x)\,dx\right)^{1/2}.
\label{eq-function-inner-product}
\end{equation}
Это скалярное произведение и порождённая им норма на \(C([a,b])\).

\end{definition}

Действительно, симметрия и билинейность следуют из свойств интеграла. Если непрерывная функция \(f\) ненулевая хотя бы в одной точке, то \(f^2\) положительна на некотором участке отрезка, поэтому \((f,f)>0\).

\begin{definition}[Ортогональная система]\label{def-orthogonal-system}

Система ненулевых функций \(\varphi_1,\varphi_2,\ldots\in C([a,b])\) называется ортогональной на \([a,b]\), если
\[
(\varphi_n,\varphi_m)=\int_a^b\varphi_n(x)\varphi_m(x)\,dx=0
\qquad(n\ne m).
\]
Для каждого \(n\) имеем \(\|\varphi_n\|^2>0\).

\end{definition}

\begin{example}[Тригонометрическая система]\label{exm-trigonometric-system}

Система
\[
1,\ \sin x,\ \cos x,\ \sin 2x,\ \cos 2x,\ldots
\]
ортогональна на \([-\pi,\pi]\). Для \(n,m\ge1\) справедливы равенства
\begin{equation}
\begin{aligned}
\int_{-\pi}^{\pi}\sin nx\cos mx\,dx&=0,\\
\int_{-\pi}^{\pi}\cos nx\cos mx\,dx
&=\begin{cases}0,&n\ne m,\\ \pi,&n=m,\end{cases}\\
\int_{-\pi}^{\pi}\sin nx\sin mx\,dx
&=\begin{cases}0,&n\ne m,\\ \pi,&n=m.\end{cases}
\end{aligned}
\label{eq-trigonometric-orthogonality}
\end{equation}
Кроме того,
\[
\int_{-\pi}^{\pi}\cos nx\,dx
=\int_{-\pi}^{\pi}\sin nx\,dx=0,
\qquad
\|1\|^2=2\pi,\quad
\|\cos nx\|^2=\|\sin nx\|^2=\pi.
\]

\end{example}

\begin{proof}

Произведение \(\sin nx\cos mx\) нечётно, поэтому его интеграл на симметричном отрезке равен нулю. Используем формулы
\[
\begin{aligned}
\cos nx\cos mx&=\tfrac12\bigl(\cos(n-m)x+\cos(n+m)x\bigr),\\
\sin nx\sin mx&=\tfrac12\bigl(\cos(n-m)x-\cos(n+m)x\bigr).
\end{aligned}
\]
Интеграл \(\cos kx\) на \([-\pi,\pi]\) равен нулю при целом \(k\ne0\) и равен \(2\pi\) при \(k=0\). Это даёт все равенства \eqref{eq-trigonometric-orthogonality} и значения норм. Ортогональность постоянной функции остальным следует из интегралов синуса и косинуса.\qedhere

\end{proof}

\begin{example}[Многочлены Лежандра]\label{exm-legendre-system}

Другой пример ортогональной системы на \([-1,1]\) --- многочлены Лежандра:
\begin{equation}
L_n(x)=\frac1{2^n n!}\frac{d^n}{dx^n}(x^2-1)^n,
\qquad n=0,1,2,\ldots
\label{eq-legendre-rodrigues}
\end{equation}
Это формула Родрига. Например, \(L_0(x)=1\), \(L_1(x)=x\), \(L_2(x)=(3x^2-1)/2\).

\end{example}
